\input{../header.tex}

\setlength{\parindent}{0em}

\begin{document}

\subsection{Implementation}
For the implementation of the \textit{Normalizing Flow} model of the galah4-dataset, the CNN architecture from the previous exercise is slightly modified.
The CNN increases the amount of feature maps from 1 \rightarrow 10 \rightarrow 20 \rightarrow 40 \rightarrow 80 in order to achieve a deeper network, which is able to learn more complex structures in the spectrum.
After that, the features maps are compressed from a size of 80 \rightarrow 40 \rightarrow 10 \rightarrow 12 \rightarrow 10, resulting in a more compact representation of the network.
Since we are using PDFs with many more parameters (up to 204 parameters for the \textit{ggt} flow function), an additional block of with a higher amount of feature maps (80 channels) was added in order to make the CNN deeper for more parameters.
In the end, the \textit{Linear}-layers are converting the features into the six \textit{Output}-neurons, representing the three labels ($T_\text{eff}, \log(g)$ and $\text{fe\_h}$) and their respective uncertainty.

For the dataset and training I used the following parameters \textit{num\_epochs} = 100 (which were not used due to early stopping), \textit{learning\_rate} = $5e-4$ and \textit{batch\_size} = 32.

\subsection{Results}
After 52 epochs, the early stopping triggered, resulting in losses of:
\begin{align*}
  \text{best\_train\_loss} &= -3.9435 (\text{in epoch 42/52})\\
  \text{best\_val\_loss} &= -4.5754 (\text{in epoch 42/52})\\
  \text{test\_loss} &= -3.1246
\end{align*}
The training loss is constantly decreasing, while the validation loss is globally decreasing, but locally oscillating. This behavior is normal and may be due to stochastic uncertainties.
Since normalizing flows are still trained via the NLL, negative loss values are possible.
For the simple \textit{t} flow function, the loss effectively resembles the Gaussian NLL with
\begin{equation}
  \nonumber
  \frac{1}{2} \left(\frac{y_i - \mu_i}{\sigma_i}\right)^2 + \log{\sigma_i},
\end{equation}
Therefore, the loss values are as expected and show a good training on a CNN network with the use of early stopping.

To validate the predicted uncertainties, a pull distribution is performed. For the different labels the following mean values $\mu$ and standard deviations $\sigma$ were calculated:
\begin{align*}
  T_\text{eff} &: \mu = -0.481 &  \sigma = 0.988\\
  \log(g) &: \mu = -0.040 & \sigma = 1.317\\
  \text{fe\_h} &: \mu = -1.705 & \sigma = 0.948
\end{align*}
In contrast to the CNN with the NLL from the last exercise:
\begin{align*}
  T_\text{eff} &: \mu = -0.103 &  \sigma = 0.686\\
  \log(g) &: \mu = -0.208 & \sigma = 0.822\\
  \text{fe\_h} &: \mu = -0.466 & \sigma = 0.497
\end{align*}
a slight improvement is visible, especially considering the standard deviations. Since the \textit{t} flow function and the NLL from last exercise should be the same, 
the improvements are most likely due to the bigger CNN.

Since I had some problems understanding the output(-dimensions) of the jammy\_flow functions, I was not able to plot the used PDF for some data samples.

\subsection{Problems and Improvements}
The biggest problem by far is the numerical instability of the optimizations for the \textit{gt}- and \textit{ggt}-flow functions.
Directly after the beginning of the training, the values are instantly \textit{NaN}.
In order to prevent this from happening, many different attempts were carried out to solve this problem:
\begin{itemize}
  \item Removing the \textit{BatchNormalization} steps from the CNN.\\
  For very small standard deviations, the \textit{BatchNormalization} can produce extremely large activations, making the exponential and logarithmic functions numerically unstabile (especially for deeper flow functions \textit{ggt > gt > t})
  \item Adding a \textit{Softplus} function to the standard deviation values and a \textit{tanh} function to the other parameters.\\
  The \textit{softplus} function
  \begin{equation}
    \nonumber
    \text{softplus}(x) = \log(1 + \exp(x))
  \end{equation}
  prevents the standard deviation from becoming negative, resulting in a \textit{NaN} for the $\log{\sigma}$.
  The \textit{tanh} function on the other side scales the parameters to the interval $[-1, 1]$, which prevents the parameters from overflowing within the exponential function, if they become too large.
  \item Adding the \textit{torch.clamp()} to the output of the CNN to prevent values smaller or equal to zero. (The same idea as the \textit{softplus} function for the standard deviations, but harder)
  \item Changing the learning rate to smaller values in order to prevent overflowing the weights.
\end{itemize}
Since all these attempts did not fix the problem at all, I tried two additional methods.

My first idea was to implement a 'warm-up' flow function, so that the more complex flow functions start with kind of pre-trained parameters.
But since the CNN is outputting different amounts of parameters depending on the flow function, the idea of pre-training the \textit{gt} flow function with the \textit{t} flow function was discarded.

The last step was to change the entire CNN structure. For that, I took the Tiny CNN from the template, but even this did not help to solve the instability problem.
\end{document}
